% ============================================================
% DOCUMENTCLASS
% ============================================================
\documentclass[12pt,oneside]{book}

% ============================================================
% METADATOS
% ============================================================
\newcommand{\universidad}{Universidad de Buenos Aires}
\newcommand{\facultad}{Facultad de Derecho}
\newcommand{\materia}{Derecho Procesal Civil y Comercial}
\newcommand{\titulo}{Partes y Sujetos en el Proceso de Desalojo}
\newcommand{\autor}{Isaac Edgar Camacho Ocampo}
\newcommand{\anio}{2026}

% ============================================================
% PAQUETES
% ============================================================
\usepackage[utf8]{inputenc}
\usepackage[T1]{fontenc}
\usepackage[spanish,es-nodecimaldot]{babel}

\usepackage[
    top=2.5cm,
    bottom=2.5cm,
    left=3cm,
    right=2.5cm,
    headheight=15pt
]{geometry}

\usepackage{amsmath}
\usepackage{amssymb}
\usepackage{mathtools}

\usepackage{graphicx}
\usepackage{float}
\usepackage{caption}
\usepackage{subcaption}

\usepackage{booktabs}
\usepackage{array}
\usepackage{longtable}
\usepackage{tabularx}

\usepackage{enumitem}

\usepackage{xcolor}
\usepackage[most]{tcolorbox}

\usepackage{fancyhdr}

\usepackage{ragged2e}

\usepackage[
    colorlinks=true,
    linkcolor=blue!50!black,
    citecolor=green!40!black,
    urlcolor=blue!60!black,
    pdftitle={Partes y Sujetos en el Proceso de Desalojo},
    pdfauthor={Isaac Edgar Camacho Ocampo},
    pdfsubject={Apuntes universitarios},
    bookmarks=true
]{hyperref}

% ============================================================
% CONFIGURACIÓN
% ============================================================
\graphicspath{
    {./img/}
    {./graficos/}
    {./figuras/}
}

\setcounter{tocdepth}{2}

\setlength{\parindent}{0pt}
\setlength{\parskip}{0.5em}

\setlist[itemize]{
    topsep=0.3em,
    itemsep=0.2em
}

\setlist[enumerate]{
    topsep=0.3em,
    itemsep=0.2em
}

\clubpenalty=10000
\widowpenalty=10000

\pagestyle{fancy}
\fancyhf{}

\fancyhead[L]{\nouppercase{\leftmark}}
\fancyhead[R]{\materia}

\fancyfoot[C]{\thepage}

\renewcommand{\headrulewidth}{0.4pt}
\renewcommand{\footrulewidth}{0pt}

\fancypagestyle{plain}{
    \fancyhf{}
    \fancyfoot[C]{\thepage}
    \renewcommand{\headrulewidth}{0pt}
}

% ============================================================
% COMANDOS
% ============================================================
\newcommand{\abs}[1]{\left|#1\right|}
\newcommand{\norm}[1]{\left\lVert#1\right\rVert}
\newcommand{\vect}[1]{\mathbf{#1}}
\newcommand{\deriv}[2]{\frac{d#1}{d#2}}
\newcommand{\pderiv}[2]{\frac{\partial#1}{\partial#2}}

% ============================================================
% ENTORNOS / CAJAS
% ============================================================
\newtcolorbox{definition}[1]{
    enhanced,
    breakable,
    title={Definición: #1},
    fonttitle=\bfseries,
    colback=blue!3,
    colframe=blue!50!black,
    boxrule=0.8pt,
    arc=2mm
}

\newtcolorbox{important}{
    enhanced,
    breakable,
    title={Importante},
    fonttitle=\bfseries,
    colback=yellow!8,
    colframe=orange!70!black,
    boxrule=0.8pt,
    arc=2mm
}

\newtcolorbox{example}[1]{
    enhanced,
    breakable,
    title={Ejemplo: #1},
    fonttitle=\bfseries,
    colback=green!4,
    colframe=green!40!black,
    boxrule=0.8pt,
    arc=2mm
}

\newtcolorbox{formula}[1]{
    enhanced,
    breakable,
    title={Fórmula / Regla: #1},
    fonttitle=\bfseries,
    colback=purple!4,
    colframe=purple!50!black,
    boxrule=0.8pt,
    arc=2mm
}

\newtcolorbox{exercise}[1]{
    enhanced,
    breakable,
    title={Ejercicio: #1},
    fonttitle=\bfseries,
    colback=gray!5,
    colframe=gray!60!black,
    boxrule=0.8pt,
    arc=2mm
}

\newtcolorbox{note}{
    enhanced,
    breakable,
    title={Nota},
    fonttitle=\bfseries,
    colback=white,
    colframe=gray!50,
    boxrule=0.6pt,
    arc=2mm
}

% ============================================================
% CONFIGURACIÓN FINAL
% ============================================================
\hypersetup{
    pdftitle={\titulo},
    pdfauthor={\autor}
}

% ============================================================
% DOCUMENT
% ============================================================
\begin{document}

% ------------------------------------------------------------
% PORTADA
% ------------------------------------------------------------
\begin{titlepage}

\thispagestyle{empty}

\begin{center}

\vspace*{1cm}

{\large \textsc{\universidad}}

\vspace{0.2cm}

{\large \textsc{\facultad}}

\vspace{3cm}

{\Huge \textbf{\materia}}

\vspace{1cm}

{\LARGE \titulo}

\vspace{0.8cm}

\textit{Comprender con precisión a cada sujeto del proceso es el primer paso para una defensa técnica sólida.}

\vspace{1.2cm}

\rule{0.70\textwidth}{0.5pt}

\vspace{0.5cm}

{\large Apuntes de estudio}

\vspace{0.5cm}

\rule{0.70\textwidth}{0.5pt}

\vfill

{\large \autor}

\vspace{0.3cm}

{\large \anio}

\end{center}

\vfill

\begin{flushleft}
\textit{Este es un modesto aporte a los estudiantes de ``Derecho Procesal Civil y Comercial'' no pretende sustituir el material oficial de la materia.}
\end{flushleft}

\end{titlepage}

% ------------------------------------------------------------
% ÍNDICE
% ------------------------------------------------------------
\tableofcontents

% ------------------------------------------------------------
% CONTENIDO
% ------------------------------------------------------------

\chapter{Sujetos del proceso de desalojo}

\section{Partes y sujetos procesales en el desalojo}

El estudio de las partes en el proceso civil, ya abordado en unidades anteriores de la materia, se aplica en este apartado puntualmente al juicio de desalojo. La pregunta que organiza el desarrollo es la siguiente: cuando se habla de ``partes'' en un desalojo, ¿se hace referencia únicamente al actor, al demandado y al juez o jueza, o existen otros sujetos que también deben ser tenidos en cuenta en esta materia?

Existe una tendencia a concebir el proceso de manera reducida, limitándolo a la tríada actor--demandado--juez. Frente a esa simplificación, corresponde manejar un esquema más amplio, en el que aparecen categorías que atraviesan todos los procesos de desalojo y de ejecución de alquileres: los \textbf{subinquilinos} y los \textbf{ocupantes}.

\begin{definition}{Subinquilinos y ocupantes}
Aunque el proceso identifique con precisión a la persona contra quien se dirige la demanda (por ejemplo, la locataria), siempre debe notificarse también a los supuestos subinquilinos y a los supuestos ocupantes del inmueble, aun cuando no estén individualizados. Se trata de una previsión constante en todo proceso de desalojo, orientada a resguardar la debida defensa en juicio de cualquier persona que pudiera estar ocupando el inmueble.
\end{definition}

A partir de esa premisa, los sujetos del proceso de desalojo pueden organizarse en tres grandes categorías:

\begin{enumerate}
    \item los órganos de la jurisdicción;
    \item los auxiliares del juez;
    \item las partes y sus colaboradores.
\end{enumerate}

Cada una de estas categorías se desarrolla en las secciones siguientes.

\section{Órganos de la jurisdicción y auxiliares del juez}

\subsection{El juez o jueza: deberes y facultades}

El primer sujeto dentro de los órganos de la jurisdicción es el juez o la jueza, con sus deberes y facultades. Esas atribuciones se ubican en el Código Procesal Civil y Comercial, particularmente en los artículos 34, 35 y 36, que regulan la denominada conciliación intraprocesal.

\begin{important}
Los artículos 34 y 35 del Código Procesal Civil y Comercial regulan los deberes y facultades del juez, incluyendo su facultad de disponer, en cualquier momento del proceso, la comparecencia personal de las partes para intentar una conciliación. El artículo 36 --- la denominada ``audiencia del artículo 36'' --- habilita al juez a convocar esta audiencia en cualquier etapa del proceso, a petición de parte o de oficio, a diferencia de otras audiencias que solo pueden fijarse en momentos procesales determinados.
\end{important}

Esta conciliación intraprocesal se da muy pocas veces en el desalojo. La razón es de tipo práctico: cuando la parte demandada solicita esta audiencia, en la generalidad de los casos el juzgado no la concede, porque el instituto suele ser utilizado como herramienta dilatoria antes que con fines genuinamente conciliatorios, en mayor medida que en otros procesos de conocimiento.

En la práctica profesional, una negociación genuina entre abogados suele darse por fuera del tribunal: los letrados coordinan directamente entre sí (por videollamada o en algún estudio jurídico) para intentar acercar posiciones, combinando agendas de manera inmediata. Ese tipo de acuerdo no depende del tribunal en el desalojo.

\subsection{Los auxiliares del juez}

Los auxiliares del juez son aquellos sujetos que colaboran con el órgano jurisdiccional para llevar adelante el proceso.

\textbf{Ministerio Público.} Dentro de esta institución existen distintas ramas; interesa aquí especialmente el Ministerio Público de la Defensa, que interviene cuando están presentes personas en situación de vulnerabilidad: niños, niñas y adolescentes, personas con capacidad restringida o adultos mayores. Cuando alguno de estos sujetos está involucrado en el desalojo, corresponde correr vista al Ministerio Público para que emita su dictamen, y recién luego el expediente vuelve al juez o jueza para resolver.

\textbf{Oficial notificador.} Es el auxiliar encargado de llevar la notificación de una resolución judicial. Un punto especialmente relevante es su competencia territorial: el oficial notificador debe pertenecer al departamento judicial correspondiente al domicilio donde se practica la notificación. Si la persona a notificar reside, por ejemplo, en un partido de la provincia de Buenos Aires distinto de la Ciudad Autónoma, no corresponde utilizar al oficial notificador de la Ciudad, sino tramitar la notificación mediante la cédula regulada por la ley 22.172 --- conocida como ``cédula ley'' --- dirigida al oficial notificador del departamento judicial correspondiente.

\begin{important}
La ley 22.172 (convenio de comunicaciones entre tribunales) regula la comunicación entre tribunales de distintas jurisdicciones mediante la denominada ``cédula ley'', que se utiliza cuando la persona a notificar tiene su domicilio en un departamento judicial distinto de aquel donde tramita el proceso.
\end{important}

Para diligenciar esa notificación existen distintas alternativas prácticas: puede llevarla el propio letrado, alguien del estudio jurídico, o puede contratarse a un gestor que cobra un arancel por el diligenciamiento; en muchos casos es el propio cliente quien asume ese costo.

\begin{example}{Falta de legitimación activa del fiador}
En un caso llevado a la práctica, una demanda de ejecución de alquileres fue promovida por el fiador --- y no por el locador ---, luego de que este hubiera pagado la deuda y pretendiera accionar contra la locataria. El supuesto plantea la discusión sobre si correspondía oponer la excepción de falta de legitimación activa para obrar, en la medida en que quien inicia la ejecución no reviste, en principio, el carácter de acreedor directo de la relación locativa.
\end{example}

% TODO: contenido incompleto o ambiguo en las notas originales: el desenlace concreto de la excepción de falta de legitimación activa planteada en el caso del fiador no queda resuelto en el apunte.

Cabe distinguir, en todo caso, entre lo estrictamente técnico-jurídico --- lo que se aprende a partir de la ley, la doctrina y la jurisprudencia --- y la realidad del ejercicio profesional, en la que un planteo jurídicamente débil puede ser igualmente utilizado con fines dilatorios: si el tribunal lo admite y la contraparte no lo cuestiona a tiempo, el acto queda firme. Esta demora estructural de los procesos judiciales ha motivado incluso llamados de atención de la Corte Interamericana de Derechos Humanos hacia la Argentina en materia de plazo razonable.

\subsection{El oficial de justicia y el lanzamiento}

El oficial de justicia se diferencia con claridad del oficial notificador. Su tarea principal en el desalojo consiste en cumplir con la diligencia del lanzamiento: la desocupación efectiva del inmueble respecto de quienes lo ocupan, sean demandados, subinquilinos u ocupantes. Se trata de una figura contratada por el poder judicial local de cada departamento judicial, que actúa por orden del juez o jueza --- nunca interviene el juez en persona ni el secretario del juzgado --- y que puede requerir, según lo peticionado y lo ordenado judicialmente, el auxilio de la fuerza pública, de un cerrajero, de un servicio médico de ambulancia o incluso de un servicio veterinario cuando hay animales en el inmueble.

\begin{table}[H]
\centering
\caption{Oficial notificador y oficial de justicia}
\begin{tabularx}{\textwidth}{>{\raggedright\arraybackslash}p{0.22\textwidth} >{\raggedright\arraybackslash}X >{\raggedright\arraybackslash}X}
\toprule
\textbf{Criterio} & \textbf{Oficial notificador} & \textbf{Oficial de justicia} \\
\midrule
Función principal & Lleva la notificación de resoluciones judiciales (cédulas) & Ejecuta el lanzamiento: hace efectiva la desocupación del inmueble \\
Competencia territorial & Departamento judicial del domicilio a notificar & Departamento judicial donde tramita el lanzamiento \\
Puede requerir & --- & Fuerza pública, cerrajero, ambulancia, servicio veterinario (si se peticiona y el juez lo ordena) \\
Relación con el flete & No interviene & Actúa junto con el fletero, contratado aparte \\
\bottomrule
\end{tabularx}
\end{table}

Respecto del flete, corresponde aclarar una confusión frecuente: no se trata de una condición procesal, sino de una persona --- el fletero --- contratada aparte del oficial de justicia y del cerrajero, necesaria porque, al abrir la puerta del inmueble, los bienes que allí se encuentren no pueden simplemente amontonarse junto a la puerta, sino que deben trasladarse a un depósito, notificando al ocupante, con un plazo mínimo de un mes, que sus bienes se encuentran allí; de no presentarse a retirarlos, se dispone su entrega a una institución de bien público.

\begin{note}
Como estrategia práctica frente al lanzamiento, conviene requerir la totalidad de los auxiliares disponibles (cerrajero, fletero, fuerza pública), dado que en muchas ocasiones no puede preverse con certeza la situación que se encontrará dentro del inmueble. Debe explicarse al cliente que, de no contratarse estos auxiliares, el lanzamiento puede no prosperar en la fecha fijada y reprogramarse varios meses después, con el consiguiente costo adicional.
\end{note}

La verificación previa de quiénes ocupan efectivamente el inmueble puede revelar, incluso, la presencia de menores de edad, supuesto que se desarrolla con un caso concreto en la Unidad 5 de este documento.

\section{Colaboradores del proceso}

En el eje de partes y colaboradores corresponde recordar una distinción ya estudiada en la parte general de la materia: el abogado puede actuar como patrocinante letrado, siempre con la firma de la parte, o como apoderado.

\begin{important}
Si un escrito se presenta bajo patrocinio letrado pero sin la firma de la parte, ese escrito será tenido como no presentado y el juzgado lo desglosará.
\end{important}

En cuanto al apoderamiento, puede tratarse de un apoderado general judicial o de un apoderado especial para el proceso de desalojo en particular. El marco normativo del poder judicial contrasta dos momentos:

\begin{important}
\textbf{Artículo 47 del Código Procesal Civil y Comercial de la Nación:} exige escritura pública (``escritura poder'') para el otorgamiento del poder judicial, conforme al régimen procesal clásico.

\textbf{Artículo 1017 del Código Civil y Comercial de la Nación} (vigente desde agosto de 2015): flexibiliza esa exigencia y admite que el mandato se otorgue bajo cualquier forma, sin requerir necesariamente escritura pública.
\end{important}

Pese a esa flexibilización sustancial, en la práctica procesal suele mantenerse la forma más rigurosa de la escritura pública para el poder judicial, por tratarse de una exigencia propia del fuero procesal.

Finalmente, corresponde referirse a los consultores técnicos, figura distinta del perito de oficio y también llamada perito de parte. Su utilidad resulta especialmente relevante cuando la liquidación del canon locativo --- que suma alquileres, intereses, servicios e impuestos --- alcanza montos elevados y comienza a controvertirse: en esos casos conviene acompañarse de un consultor técnico, con formación en matemáticas o en el manejo de herramientas de análisis, para no afrontar en soledad esa tarea.

% ------------------------------------------------------------
% CORNELL CAPÍTULO 1
% ------------------------------------------------------------
\section*{Cuadro Cornell --- Sujetos del proceso}
\addcontentsline{toc}{section}{Cuadro Cornell --- Sujetos del proceso}

\begin{table}[H]
\centering
\begin{tabularx}{\textwidth}{
    >{\raggedright\arraybackslash}p{0.30\textwidth}
    >{\raggedright\arraybackslash}X
}
\toprule
\textbf{Preguntas / palabras clave} & \textbf{Notas / respuestas} \\
\midrule

\textbf{¿Quiénes son ``partes'' en un desalojo, además de actor, demandado y juez?} &
También hay órganos de la jurisdicción (juez, con deberes y facultades de los arts. 34, 35 y 36 CPCCN), auxiliares (Ministerio Público, oficial notificador, oficial de justicia) y, como categoría transversal, subinquilinos y ocupantes, que deben notificarse aunque no estén individualizados. \\

\textbf{¿Qué es la audiencia del artículo 36 y por qué casi no se concede en el desalojo?} &
Es una audiencia conciliatoria que el juez puede convocar en cualquier momento del proceso, de oficio o a pedido de parte. En el desalojo suele denegarse porque en la práctica se interpreta que se solicita con fines dilatorios más que conciliatorios. \\

\textbf{¿Qué diferencia al oficial notificador del oficial de justicia?} &
El notificador lleva cédulas de resoluciones judiciales y depende del departamento judicial del domicilio a notificar. El oficial de justicia ejecuta el lanzamiento (la desocupación efectiva) y puede requerir fuerza pública, cerrajero, ambulancia o veterinario si se peticiona y el juez lo ordena. \\

\textbf{¿Qué es la cédula ley y cuándo se usa?} &
Es la notificación regulada por la ley 22.172, utilizada cuando la persona a notificar tiene su domicilio en un departamento judicial distinto de aquel donde tramita el proceso. \\

\textbf{¿Qué diferencia hay entre patrocinio letrado y apoderado?} &
El patrocinio requiere siempre la firma de la parte junto a la del letrado; sin ella, el escrito se tiene por no presentado. El apoderado actúa representando directamente a la parte, con poder otorgado conforme al art. 47 CPCCN (escritura pública) o, con mayor flexibilidad, al art. 1017 CCCN (cualquier forma) desde 2015. \\

\textbf{¿Para qué sirve un consultor técnico en un desalojo o ejecución de alquileres?} &
Para controvertir liquidaciones complejas de canon locativo, intereses y servicios cuando los montos son elevados y se discuten; aporta una mirada técnica distinta de la del perito de oficio. \\

\midrule
\multicolumn{2}{p{\textwidth}}{
\textbf{Resumen:} el proceso de desalojo involucra, además de actor, demandado y juez, a órganos de la jurisdicción con deberes y facultades específicas, a auxiliares (Ministerio Público, oficial notificador, oficial de justicia) con funciones diferenciadas, y a colaboradores de las partes (patrocinio letrado, apoderamiento, consultores técnicos), cada uno regido por reglas propias de competencia, forma y oportunidad.
} \\
\bottomrule
\end{tabularx}
\end{table}

\chapter{Competencia y caso práctico}

\section{Competencia en el juicio de desalojo}

El primer paso para determinar la competencia en un desalojo consiste siempre en revisar el contrato de locación completo, si existe. Por tratarse de una cuestión patrimonial, las partes pueden haber pactado un \textbf{pacto de foro prorrogado}, es decir, una cláusula que fija de antemano el tribunal competente para dirimir cualquier conflicto, o bien una \textbf{cláusula compromisoria} que derive el conflicto a un tribunal arbitral.

Cuando no hay contrato disponible se aplican las reglas generales de competencia.

\begin{important}
El artículo 5 y siguientes del Código Procesal Civil y Comercial establecen que, a falta de pacto de foro prorrogado o cláusula compromisoria, en cuestiones patrimoniales como el desalojo la competencia se determina por el lugar de ubicación del inmueble o, en su defecto, por el domicilio del demandado.
\end{important}

\begin{definition}{Prórroga tácita de competencia}
Si, pese a existir un pacto de foro prorrogado, el actor inicia la demanda ante un tribunal distinto del pactado y la parte demandada contesta la demanda sin oponer excepción de incompetencia, se produce una prórroga tácita que consolida la competencia de ese tribunal. La alternativa consiste en oponer la excepción de incompetencia invocando la cláusula compromisoria o el pacto de foro prorrogado.
\end{definition}

La competencia territorial del oficial notificador depende del domicilio real que en ese momento tenga el sujeto a notificar --- el domicilio que se denuncie como tal o el que se corrobore mediante el Registro Nacional de las Personas ---, y no necesariamente coincide con el domicilio donde tramita el proceso. En el desalojo, muchas veces también deben notificarse los fiadores o garantes, en su carácter de obligados como principales pagadores.

\section{Intervención del Ministerio Público de la Defensa: un caso real}

La verificación de quiénes ocupan efectivamente el inmueble, antes del lanzamiento, resulta determinante para el curso del proceso, según ilustra el siguiente caso.

\begin{example}{Desalojo con menores en el inmueble (Ciudadela, partido de Tres de Febrero)}
Un proceso de desalojo tramitaba en el fuero nacional civil respecto de un inmueble ubicado en Ciudadela, partido de Tres de Febrero, provincia de Buenos Aires. Al llegar la etapa del lanzamiento, intervino el oficial de justicia correspondiente al departamento judicial de Tres de Febrero, quien, en el recorrido previo al inmueble realizado antes de fijar fecha, advirtió la presencia de menores de edad en la vivienda.

Como el proceso tramitaba en el fuero nacional civil pero el inmueble se encontraba en jurisdicción provincial, la solución procesal consistió en abrir un oficio bajo la ley 22.172 dirigido al juzgado en lo civil y comercial de San Martín, para que habilitara la intervención y la consulta al Ministerio Público de la Defensa por la presencia de menores.

El trasfondo del caso correspondía a un locador de edad avanzada, viudo y sin hijos, que vivía de su jubilación y del ingreso del alquiler.
\end{example}

\begin{important}
La verificación de quiénes ocupan efectivamente un inmueble --- incluidos menores, personas con capacidad restringida o adultos mayores --- no es un trámite formal secundario: puede modificar por completo el curso del proceso, obligando a articular mecanismos interjurisdiccionales (como la ley 22.172) y a dar intervención al Ministerio Público de la Defensa antes de poder concretar el lanzamiento.
\end{important}

% ------------------------------------------------------------
% CORNELL CAPÍTULO 2
% ------------------------------------------------------------
\section*{Cuadro Cornell --- Competencia y caso práctico}
\addcontentsline{toc}{section}{Cuadro Cornell --- Competencia y caso práctico}

\begin{table}[H]
\centering
\begin{tabularx}{\textwidth}{
    >{\raggedright\arraybackslash}p{0.30\textwidth}
    >{\raggedright\arraybackslash}X
}
\toprule
\textbf{Preguntas / palabras clave} & \textbf{Notas / respuestas} \\
\midrule

\textbf{¿Cómo se determina la competencia en un desalojo?} &
Primero se revisa el contrato de locación por si existe un pacto de foro prorrogado o una cláusula compromisoria. A falta de eso, rigen los arts. 5 y siguientes del CPCCN: competencia por el lugar del inmueble o, subsidiariamente, por el domicilio del demandado. Puede haber prórroga tácita si el demandado contesta sin oponer incompetencia. \\

\textbf{¿Por qué puede intervenir el Ministerio Público de la Defensa en un lanzamiento?} &
Porque el oficial de justicia, al hacer el recorrido previo, puede advertir la presencia de menores, personas con capacidad restringida o adultos mayores en el inmueble. Si el inmueble está en otra jurisdicción, puede requerir un oficio bajo la ley 22.172 para habilitar esa intervención. \\

\midrule
\multicolumn{2}{p{\textwidth}}{
\textbf{Resumen:} la competencia en el desalojo se determina, en primer lugar, por lo pactado en el contrato (foro prorrogado o cláusula compromisoria) y, en su defecto, por las reglas generales del lugar del inmueble o del domicilio del demandado. La verificación de las personas que efectivamente ocupan el inmueble puede motivar la intervención del Ministerio Público de la Defensa y, si median distintas jurisdicciones, el uso de mecanismos interjurisdiccionales como la ley 22.172.
} \\
\bottomrule
\end{tabularx}
\end{table}

\chapter{Calidad, legitimidad y domicilios}

\section{Subinquilinos, ocupantes y terceros: el fiador}

La figura de subinquilinos y ocupantes está siempre presente en la demanda de desalojo. En el escrito de demanda es habitual incluir una fórmula genérica que da cuenta de que se desconoce si existen subinquilinos u ocupantes distintos del demandado principal, y que, en caso de conocerse, se los denuncia para su identificación.

\begin{definition}{Fórmula genérica ``subinquilinos y ocupantes''}
La razón de admitir esta imprecisión genérica es procesal: identificar a cada persona por nombre y apellido obligaría a notificar correctamente a cada una de ellas, y omitir a alguna comprometería la debida defensa en juicio. Por ello, no corresponde individualizar a cada conviviente; se demanda genéricamente contra el locatario y contra sus subinquilinos y ocupantes.
\end{definition}

En cuanto a los terceros, corresponde mencionar a peritos y testigos, y detenerse en la figura del fiador o garante.

\begin{important}
No existe una obligación legal de notificar al fiador del inicio del proceso de desalojo; se trata de una \textbf{buena práctica} que algunos juzgados adoptan según el criterio de cada tribunal. El ``debe'' es una obligación cuya inobservancia acarrea sanción; la notificación al fiador, en cambio, no está impuesta bajo sanción, sino que constituye una práctica que algunos tribunales adoptan porque resulta útil: apenas notificado formalmente, el fiador suele intentar algún entendimiento antes de afrontar una condena.
\end{important}

Con independencia de la vía judicial, suele intentarse contactar a la contraparte por todos los medios disponibles antes de iniciar el proceso, con el fin de negociar; sin embargo, la notificación judicial formal al fiador tiende a generar una reacción distinta de la que producen los intentos informales de acercamiento.

Respecto de la posibilidad de solicitar una medida cautelar sobre el inmueble de garantía del fiador, corresponde señalar que las medidas cautelares constituyen un subsistema autónomo del sistema procesal y requieren siempre reunir sus propios presupuestos --- verosimilitud del derecho y peligro en la demora ---, con independencia del proceso principal de que se trate.

\begin{table}[H]
\centering
\caption{Medidas cautelares sobre bienes del fiador según el tipo de proceso}
\begin{tabularx}{\textwidth}{>{\raggedright\arraybackslash}p{0.30\textwidth} X}
\toprule
\textbf{Proceso} & \textbf{¿Procede cautelar sobre bienes del fiador?} \\
\midrule
Desalojo & No corresponde: lo que se persigue es recuperar la tenencia del inmueble, no una prestación dineraria del fiador. \\
Ejecución de alquileres & Sí puede jugar el subsistema de medidas cautelares, aunque en la generalidad de los tribunales no se conceden hasta que la litis está trabada, es decir, hasta que la contraparte esté debidamente notificada, haya contestado o no. \\
\bottomrule
\end{tabularx}
\end{table}

Como alternativas dentro del universo de medidas cautelares se mencionan la inhibición general de bienes (cuando no hay bienes identificados del deudor) y, más excepcionalmente, la prohibición de contratar, útil por ejemplo si se advierte que el fiador está por enajenar el inmueble que constituye su garantía patrimonial.

Estos procesos pueden involucrar tanto a personas físicas como a personas jurídicas o asociaciones, y la perspectiva de vulnerabilidad, junto con las Reglas de Brasilia, se mantiene como marco transversal en esta materia.

\section{Calidad de parte y legitimidad}

\begin{definition}{Calidad de parte y legitimidad}
Una cosa es la calidad de parte y otra es la legitimidad de parte. Un menor de edad no tiene calidad de parte en el proceso porque no puede actuar por sí mismo debido a su incapacidad de hecho; debe estar presente su representante legal. Sin embargo, ese mismo menor puede tener legitimidad sustancial --- por ejemplo, si el inmueble en cuestión le fue donado por sus abuelos, es el titular de la relación jurídica sustancial, aunque no pueda comparecer personalmente en el proceso.
\end{definition}

\begin{important}
Un sujeto debe reunir dos condiciones: capacidad de derecho --- legitimidad sustancial, ser el titular de la relación sustancial --- y capacidad de hecho, para poder ejercer por sí mismo esa titularidad dentro del proceso. No hace falta ser el titular dominial del inmueble para promover un desalojo: puede ser suficiente ser condómino, usuario o usufructuario, en tanto se sea el titular de la relación contractual de locación --- quien efectivamente firmó el contrato como locador.
\end{important}

Si el titular de esa relación sustancial fuera menor de dieciocho años y no estuviera emancipado, carecerá de capacidad de hecho, y deberá comparecer en su representación su padre, su madre o, en su defecto, su tutor.

Vinculada a esta distinción se encuentra la cuestión de la representación. Es necesario revisar con cuidado los poderes acompañados en el expediente, porque en ocasiones se presenta un apoderamiento que no habilita la representación pretendida --- por ejemplo, un poder otorgado para un proceso laboral que nada tiene que ver con la representación convencional necesaria en un desalojo ---. Dado que el tribunal no siempre controla esta cuestión de oficio, corresponde a los abogados de cada parte verificarla y, eventualmente, plantear la falta de legitimación procesal.

Pueden distinguirse así tres formas de representación: la representación legal (padres, o el curador si se trata de una persona con capacidad restringida) y la representación convencional, propia de la relación abogado--cliente mediante apoderamiento, que puede instrumentarse por instrumento privado con firma certificada o por escritura pública, conforme a la mayor flexibilidad introducida por el Código Civil y Comercial desde 2015 (art. 1017 CCCN, ya visto en el Capítulo 1).

Finalmente, corresponde diferenciar dos excepciones procesales posibles frente a estos problemas: la \textbf{falta de legitimación procesal} (vinculada a la capacidad de hecho o a la representación insuficiente) y la \textbf{falta de legitimación sustancial}, activa o pasiva, cuando quien demanda o es demandado no es efectivamente quien firmó el contrato de locación ni el titular de la relación jurídica sustancial.

\section{Los domicilios en el proceso}

Un error frecuente consiste en confundir distintos tipos de domicilio que cumplen funciones procesales diferentes. Es habitual que en los escritos judiciales se emplee la fórmula ``constituyo domicilio legal en...'' cuando en realidad se está constituyendo un domicilio procesal o \emph{ad litem}, y no un domicilio legal en sentido técnico.

\begin{table}[H]
\centering
\caption{Tipos de domicilio en el proceso}
\begin{tabularx}{\textwidth}{>{\raggedright\arraybackslash}p{0.26\textwidth} X}
\toprule
\textbf{Tipo de domicilio} & \textbf{Función y características} \\
\midrule
Domicilio especial & Pactado en el contrato de locación, a los fines sustanciales del vínculo contractual. No equivale al domicilio procesal ni al legal. \\
Domicilio procesal o \emph{ad litem} & Constituido por las partes a los fines del proceso; debe ubicarse dentro del radio del juzgado donde tramita la causa. \\
Domicilio legal & El que fija la ley para ciertos sujetos: por ejemplo, el domicilio de una sociedad registrada, o el domicilio de una persona con capacidad restringida, que es el de su curador. \\
Domicilio electrónico constituido & Obligatorio en todo proceso judicial. En el orden nacional, el sistema es el M.E.V.; en los departamentos judiciales de la provincia de Buenos Aires, el sistema es Notificaciones CBA. \\
Domicilio real & Lugar de residencia efectiva de la persona. Se utiliza para actos específicos: la sentencia y el traslado de la demanda. \\
\bottomrule
\end{tabularx}
\end{table}

El domicilio electrónico es obligatorio en cualquier departamento judicial y reemplazó, en gran medida, a la cédula en papel: todas las notificaciones que antes se practicaban mediante el oficial notificador ahora se cursan en forma electrónica o digital, según el sistema del departamento judicial correspondiente.

\begin{important}
Solo puede notificarse por vía electrónica cuando la contraparte ya constituyó su domicilio electrónico en el proceso, es decir, cuando ya se presentó con su letrado. Si todavía no se presentó, corresponde practicar la notificación por cédula en papel al domicilio procesal constituido (que no necesariamente coincide con el estudio jurídico del abogado, aunque en la generalidad de los casos sí coincide) o, si tampoco hay domicilio constituido, al domicilio real, donde debe haber alguien que reciba la cédula, identificándose con su documento.
\end{important}

Muy pocos actos del proceso se notifican al domicilio real: entre ellos, la sentencia y el traslado de la demanda, precisamente porque en esas etapas todavía no existe domicilio procesal constituido por la contraparte. En el patrocinio jurídico gratuito es frecuente constituir domicilio en los estrados del juzgado, lo cual habilita que las notificaciones corran automáticamente los días de nota, con el riesgo de que el letrado pierda el plazo si no controla esos días con atención.

\begin{note}
Una ley anterior --- hoy derogada por el decreto 70/2023 --- establecía como obligatorio incluir en los contratos de locación un domicilio de correo electrónico, y disponía que toda notificación cursada a esa dirección servía como prueba definitoria. Después de la sanción de dicho decreto, muchos contratos dejaron de redactarse incluyendo esa cláusula.
\end{note}

Cuando las partes pactaban expresamente sus direcciones de correo electrónico en el contrato, ninguna de las dos podía alegar luego desconocimiento de la notificación, dado que la obligación de avisar cualquier cambio de dirección también quedaba a cargo de cada parte. Sin embargo, esa misma norma, pensada originalmente en beneficio del locatario --- para que pudiera notificar rápidamente al locador reclamos como una filtración de agua o un desperfecto en las instalaciones ---, terminó siendo en la práctica perjudicial para el locatario, en la medida en que habilitaba notificar por esa vía cualquier tipo de acto, incluyendo actos procesales relevantes que antes requerían notificación formal.

% TODO: contenido incompleto o ambiguo en las notas originales: el apunte no precisa el criterio definitivo para determinar la responsabilidad entre locador y locatario frente a un desperfecto edilicio; solo indica que en la práctica suele recurrirse a un profesional idóneo (por ejemplo, un plomero) para establecer la causa del daño.

% ------------------------------------------------------------
% CORNELL CAPÍTULO 3
% ------------------------------------------------------------
\section*{Cuadro Cornell --- Calidad, legitimidad y domicilios}
\addcontentsline{toc}{section}{Cuadro Cornell --- Calidad, legitimidad y domicilios}

\begin{table}[H]
\centering
\begin{tabularx}{\textwidth}{
    >{\raggedright\arraybackslash}p{0.30\textwidth}
    >{\raggedright\arraybackslash}X
}
\toprule
\textbf{Preguntas / palabras clave} & \textbf{Notas / respuestas} \\
\midrule

\textbf{¿Por qué se admite la fórmula genérica ``subinquilinos y ocupantes'' sin identificarlos por nombre?} &
Porque exigir la identificación nominal de cada conviviente vulneraría la debida defensa en juicio de personas no identificadas; se prefiere admitir esa imprecisión y notificar genéricamente a todo ocupante. \\

\textbf{¿Es obligatorio notificar al fiador del inicio del desalojo?} &
No hay una obligación legal (un ``debe'' sancionable); es una buena práctica que algunos juzgados adoptan porque suele apurar un acuerdo con el garante. \\

\textbf{¿Puede pedirse una medida cautelar sobre el inmueble del fiador en un desalojo?} &
No como parte del desalojo (que busca recuperar la tenencia). Las cautelares (inhibición general de bienes, prohibición de contratar) juegan como subsistema propio, típicamente en la ejecución de alquileres, y usualmente exigen que la litis esté trabada. \\

\textbf{¿Qué diferencia hay entre calidad de parte y legitimidad?} &
La calidad de parte exige capacidad de hecho para actuar por sí en el proceso (un menor no la tiene y necesita representante). La legitimidad exige capacidad de derecho: ser el titular de la relación jurídica sustancial, lo cual puede tenerlo incluso un menor (por ejemplo, si el inmueble le fue donado). \\

\textbf{¿Hace falta ser dueño del inmueble para promover un desalojo?} &
No. Basta con ser el titular de la relación contractual de locación (el locador que firmó el contrato), aunque no sea el titular dominial: puede ser condómino, usuario o usufructuario. \\

\textbf{¿Cuál es la diferencia entre domicilio procesal (ad litem), legal, especial y real?} &
El procesal o ad litem se constituye a los fines del proceso, dentro del radio del juzgado. El legal es el que fija la ley (por ejemplo, el de una sociedad registrada). El especial es el pactado en el contrato para notificaciones sustanciales. El real es el de residencia efectiva y solo se usa para actos específicos como la sentencia y el traslado de demanda. \\

\textbf{¿Cuándo puede notificarse electrónicamente a la contraparte?} &
Solo cuando esa parte ya constituyó su domicilio electrónico en el proceso, es decir, cuando ya se presentó con su letrado. Antes de eso, corresponde cédula en papel al domicilio procesal o, en su defecto, al domicilio real. \\

\midrule
\multicolumn{2}{p{\textwidth}}{
\textbf{Resumen:} la identificación genérica de subinquilinos y ocupantes protege la defensa en juicio de terceros no individualizados; la notificación al fiador y las cautelares sobre sus bienes no son exigencias legales automáticas, sino prácticas o subsistemas condicionados. La calidad de parte (capacidad de hecho) y la legitimidad (titularidad sustancial) son categorías distintas, y el sistema de domicilios (especial, procesal, legal, electrónico y real) responde a funciones procesales diferenciadas cuya confusión puede acarrear consecuencias prácticas relevantes.
} \\
\bottomrule
\end{tabularx}
\end{table}

\chapter{Cargas, principios y figuras procesales}

\section{Cargas, deberes y obligaciones. Principio dispositivo y preclusión}

La carga, el deber y la obligación no son sinónimos en el proceso civil.

\begin{definition}{Carga, deber y obligación}
\textbf{Carga:} imperativo del propio interés. Ejemplo: la carga de contestar la demanda en el desalojo. Al demandado le llega la demanda y puede optar por no responder (asumiendo el riesgo de no ejercer su defensa) o contestarla y defenderse: es su propio interés el que está en juego, no una sanción externa.

\textbf{Deber:} implica una obligación cuyo incumplimiento acarrea sanción.

\textbf{Obligación:} nace de una relación jurídica concreta. Por ejemplo, quien tiene una resolución de desalojo firme está obligado a entregar la tenencia dentro del plazo fijado (por ejemplo, diez días); si no lo hace, corresponde ejecutar esa obligación mediante el lanzamiento.
\end{definition}

Esta distinción se vincula con el \textbf{principio dispositivo} que rige el sistema procesal: son las partes quienes disponen del impulso del proceso, y no cabe esperar pasivamente que sea el tribunal --- o la contraparte --- quien lo impulse.

El segundo principio destacado es el sistema de \textbf{preclusión}: una vez vencido el plazo para presentar determinado planteo (por ejemplo, la excepción de falta de legitimación activa tratada en capítulos anteriores), esa posibilidad queda clausurada, aunque el tipo de proceso asignado (sumarísimo u ordinario) module la extensión de esos plazos --- por ejemplo, quince días en un proceso de conocimiento ordinario.

\begin{note}
Para el cómputo de plazos procesales y el uso de las llamadas ``dos primeras horas'' o ``cuatro primeras horas'' habilitadas para presentar escritos fuera del horario ordinario, es necesario conocer el horario de atención de los tribunales: en el orden nacional, de 7:30 a 13:30 horas; en la provincia de Buenos Aires, de 8:00 a 14:00 horas.
\end{note}

\section{Litisconsorcio y proceso colectivo}

El \textbf{litisconsorcio} es la pluralidad de partes que puede presentarse tanto del lado actor como del lado demandado, o de ambos lados a la vez. En el desalojo, esta pluralidad puede darse, por ejemplo, cuando existen varios condóminos que revisten en conjunto el carácter de locadores frente al locatario y sus ocupantes.

\begin{definition}{Litisconsorcio voluntario y necesario}
El \textbf{litisconsorcio voluntario} es la regla general: la pluralidad se da de manera facultativa. El \textbf{litisconsorcio necesario} está reservado a casos muy particulares en los que la ley exige demandar conjuntamente a todos los sujetos involucrados. Esta última modalidad no se presenta en el desalojo.
\end{definition}

\begin{example}{Litisconsorcio necesario en una acción de escrituración}
Si dos personas son condóminas de un mismo inmueble, quien pretende la escrituración debe necesariamente demandar a ambas, porque cada una solo puede escriturar su propia proporción (por ejemplo, el cincuenta por ciento).
\end{example}

El \textbf{proceso colectivo}, por su parte, no tiene aplicación directa en materia de desalojo. Fue incorporado hace más de treinta años por la reforma constitucional de 1994 y, sin embargo, no cuenta todavía con una regulación procesal operativa integral; la Corte Suprema, mediante una acordada, avanzó en la creación de un registro de procesos colectivos, pero el instituto sigue sin desarrollo legislativo pleno.

\begin{note}
Como ejemplos de proceso colectivo --- ajenos a la materia de desalojo, pero útiles para no confundir figuras --- pueden mencionarse los casos ambientales (por ejemplo, el vertido de desechos industriales en un curso de agua que contamina las napas de un barrio o una ciudad entera) y los reclamos colectivos por comisiones bancarias cobradas indebidamente en cajas de ahorro, que en ocasiones derivan en sentencias colectivas que benefician a todos los titulares de cuentas alcanzados.
\end{note}

\section{Desalojos administrativos y sus límites}

Existen desalojos realizados por el Gobierno de la Ciudad de Buenos Aires al amparo de un decreto administrativo. Ese decreto existe y es real, pero su aplicación está limitada a causales específicas y extremas: situaciones de emergencia habitacional grave, como el riesgo de derrumbe estructural.

\begin{example}{Edificio abandonado con riesgo de derrumbe}
Un edificio abandonado por la constructora que lo levantó (que quebró y lo dejó inconcluso) fue tomado luego por numerosas familias, sin servicios mínimos y con riesgo edilicio en losas y columnas. En ese tipo de situación extrema, el decreto administrativo habilita lo que se denomina ``desalojo administrativo'', sin intervención judicial previa.
\end{example}

\begin{important}
El decreto administrativo que habilita desalojar sin intervención judicial previa solo puede fundarse en causales extremas de emergencia (riesgo de derrumbe, ausencia de condiciones mínimas de habitabilidad), y no puede utilizarse frente a la mera ocupación irregular de un inmueble, situación que sí requiere la intervención del Poder Judicial y de las defensorías correspondientes.
\end{important}

\begin{example}{Medida cautelar judicial frente a un desalojo administrativo}
En un caso concreto, un juez de la Ciudad de Buenos Aires intervino al considerar que la Ciudad estaba invocando el decreto de emergencia fuera de sus causales, y dictó una medida cautelar dentro de un proceso colectivo para frenar ese desalojo administrativo hasta que la situación fuera revisada judicialmente.
\end{example}

Frente a un edificio verdaderamente en riesgo de derrumbe, con familias viviendo sin servicios mínimos, el Poder Ejecutivo debe actuar para resolver una situación de peligro real; pero el mismo mecanismo no puede emplearse como atajo frente a conflictos de ocupación que no presentan ese riesgo extremo, porque para esos casos está creado el Poder Judicial y están creadas las defensorías, precisamente para analizar la situación de cada persona involucrada, sea o no ocupante irregular.

% ------------------------------------------------------------
% CORNELL CAPÍTULO 4
% ------------------------------------------------------------
\section*{Cuadro Cornell --- Cargas, principios y figuras procesales}
\addcontentsline{toc}{section}{Cuadro Cornell --- Cargas, principios y figuras procesales}

\begin{table}[H]
\centering
\begin{tabularx}{\textwidth}{
    >{\raggedright\arraybackslash}p{0.30\textwidth}
    >{\raggedright\arraybackslash}X
}
\toprule
\textbf{Preguntas / palabras clave} & \textbf{Notas / respuestas} \\
\midrule

\textbf{¿Qué diferencia hay entre carga, deber y obligación?} &
La carga es un imperativo del propio interés (contestar la demanda). El deber implica sanción por incumplimiento. La obligación nace de una relación jurídica concreta (entregar la tenencia tras una resolución firme de desalojo). \\

\textbf{¿Cuándo hay litisconsorcio necesario y se da en el desalojo?} &
El litisconsorcio necesario exige demandar conjuntamente a todos los sujetos (ejemplo: escrituración contra todos los condóminos). En el desalojo, la regla es el litisconsorcio voluntario; el necesario no se presenta en esta materia. \\

\textbf{¿Puede el gobierno de una ciudad desalojar sin intervención judicial?} &
Solo excepcionalmente, mediante decreto administrativo limitado a causales extremas de emergencia (por ejemplo, riesgo real de derrumbe). Fuera de esas causales, corresponde la intervención del Poder Judicial y de las defensorías. \\

\midrule
\multicolumn{2}{p{\textwidth}}{
\textbf{Resumen:} el proceso de desalojo se sostiene sobre la distinción entre cargas, deberes y obligaciones, y sobre los principios dispositivo y de preclusión que rigen el impulso y los plazos procesales. El litisconsorcio voluntario es la regla en el desalojo, mientras que el litisconsorcio necesario y el proceso colectivo son figuras ajenas a esta materia. Los desalojos administrativos solo son válidos frente a causales extremas de emergencia habitacional, quedando sujetos siempre a control judicial.
} \\
\bottomrule
\end{tabularx}
\end{table}

\chapter{Fianza y principios procesales}

\section{Extinción de la fianza (artículo 1225 CCCN) y buena fe procesal}

El artículo 1225 del Código Civil y Comercial de la Nación constituye una de las normas de mayor aplicación práctica tanto en el desalojo como en la ejecución de alquileres, vinculada a la extinción de la fianza cuando el contrato de locación no se renueva formalmente.

\begin{formula}{Artículo 1225 CCCN --- Extinción de la fianza}
La responsabilidad del fiador se limita al plazo contractual pactado. Si el contrato de locación continúa más allá de ese plazo sin que se celebre un nuevo contrato o una renovación formal en la que el fiador vuelva a obligarse expresamente, la fianza no se renueva automáticamente junto con la locación de hecho.
\end{formula}

\begin{example}{Extinción de la fianza por falta de renovación formal}
Un contrato de locación por dos años, con el hermano del locatario como fiador, vence en diciembre de 2025. Como la relación entre locador y locatario es buena y este último es un buen pagador, el locador le propone continuar de palabra --- actualizando el valor del alquiler según cómo evolucione el mercado --- sin renovar por escrito ni formalizar nuevamente la fianza del hermano. La locación de hecho continúa, con pagos mensuales, pero sin la formalidad de la fianza renovada.

Si en algún momento posterior el locatario deja de pagar y el locador inicia el proceso de desalojo (para recuperar la tenencia) y, en paralelo, la ejecución de la deuda por alquileres y servicios adeudados, la posibilidad de demandar también al hermano fiador depende de si el abogado del fiador conoce y opone oportunamente la norma del artículo 1225 CCCN.
\end{example}

Si el abogado del fiador desconoce el límite del artículo 1225, es posible que termine aconsejándole pagar la deuda reclamada, o que se presente en el proceso y conteste la demanda sin oponer la excepción correspondiente, asumiendo una responsabilidad de la que en rigor podría liberarse. En cambio, si el abogado del fiador conoce la norma y ejerce oportunamente la carga de comparecer y defenderse, puede oponer la \textbf{excepción de falta de legitimación pasiva} invocando el artículo 1225, con fundamento en que la responsabilidad del fiador se limitaba al plazo contractual original y ese plazo ya había vencido sin renovación formal.

Corresponde, asimismo, un rol activo del abogado del locador frente a esta norma: si asesora a la parte interesada en demandar también al fiador, debe advertirle honestamente sobre el riesgo de que la excepción del artículo 1225 prospere, dejando a criterio del cliente si de todas formas desea afrontar la tasa de justicia y demandar igualmente al garante, con la posibilidad de que este se defienda exitosamente o no.

\begin{important}
El artículo 1225 CCCN no depende de una cláusula contractual que extienda la fianza ``hasta que el locatario devuelva la tenencia al locador'': si el locador no fue diligente en renovar el contrato por escrito ni en iniciar oportunamente el proceso de desalojo, el fiador puede liberarse de responder por el período posterior al vencimiento del plazo original pactado. Hay que conocer la norma para poder aplicarla oportunamente.
\end{important}

% TODO: contenido incompleto o ambiguo en las notas originales: el apunte anuncia el envío posterior de tres fallos vinculados a la falta de legitimación activa y al alcance del lanzamiento respecto de subinquilinos y ocupantes, sin transcribir su contenido.

El lanzamiento se dirige tanto contra el demandado identificado como contra todo tercero que ocupe el inmueble sin haber sido individualizado, conforme a lo desarrollado en el Capítulo 3 respecto de subinquilinos y ocupantes.

\subsection*{Principios procesales de cierre: colaboración, buena fe y disponibilidad de las formas}

El diseño del Código no está pensado en torno a una colaboración intensa entre las partes, sino más bien bajo una lógica de contraposición, salvo que las partes logren algún acuerdo. Aun litigando, conviene siempre intentar negociar en paralelo al desalojo judicial.

\begin{definition}{Principio de disponibilidad de las formas}
En cuestiones de carácter patrimonial como el desalojo, las partes pueden disponer del proceso, aunque esto no significa que puedan disponer libremente de todas las formas procesales, sino de aspectos puntuales acordados entre los abogados.
\end{definition}

\begin{example}{Pericial arbitral frente a un daño edilicio}
Ante una controversia sobre si un daño edilicio (por ejemplo, una rotura en las instalaciones) fue provocado por el desgaste natural o por un uso indebido del inquilino, las partes pueden acordar previamente a la etapa de prueba una pericial arbitral --- prevista en el Código Procesal --- para que un perito árbitro determine la causa del daño, evitando así un litigio más extenso sobre ese punto puntual.
\end{example}

% ------------------------------------------------------------
% CORNELL CAPÍTULO 5
% ------------------------------------------------------------
\section*{Cuadro Cornell --- Fianza y principios procesales}
\addcontentsline{toc}{section}{Cuadro Cornell --- Fianza y principios procesales}

\begin{table}[H]
\centering
\begin{tabularx}{\textwidth}{
    >{\raggedright\arraybackslash}p{0.30\textwidth}
    >{\raggedright\arraybackslash}X
}
\toprule
\textbf{Preguntas / palabras clave} & \textbf{Notas / respuestas} \\
\midrule

\textbf{¿Hasta cuándo responde el fiador si el contrato de locación no se renueva por escrito (art. 1225 CCCN)?} &
La responsabilidad del fiador se limita al plazo contractual pactado. Si la locación continúa de hecho sin renovación formal ni nueva fianza expresa, el fiador puede liberarse de responder por el período posterior, siempre que su abogado oponga oportunamente esa defensa. \\

\textbf{¿Qué es la disponibilidad de las formas y cómo se aplica en el ejemplo del daño edilicio?} &
Es la posibilidad de que las partes acuerden aspectos puntuales del trámite en cuestiones patrimoniales, por ejemplo, sustituir la etapa de prueba tradicional por una pericial arbitral que determine la causa de un daño edilicio antes de litigar sobre ese punto. \\

\midrule
\multicolumn{2}{p{\textwidth}}{
\textbf{Resumen:} el artículo 1225 CCCN limita la responsabilidad del fiador al plazo contractual original, liberándolo cuando la locación continúa de hecho sin renovación formal, siempre que se oponga oportunamente la defensa correspondiente. El proceso de desalojo, aunque estructurado sobre una lógica de contraposición entre las partes, admite --- en cuestiones patrimoniales --- la disponibilidad de ciertas formas procesales mediante acuerdo entre los abogados.
} \\
\bottomrule
\end{tabularx}
\end{table}

\chapter*{Síntesis general}
\addcontentsline{toc}{chapter}{Síntesis general}

La materia recorre, con un hilo conductor claro, el elenco de sujetos que rodean el proceso de desalojo y que exceden la tríada actor--demandado--juez: órganos de la jurisdicción con sus deberes y facultades (arts. 34, 35 y 36 CPCCN), auxiliares como el Ministerio Público, el oficial notificador y el oficial de justicia --- con sus tareas diferenciadas y los terceros que este último puede requerir para el lanzamiento: cerrajero, fletero, fuerza pública, servicios médicos o veterinarios ---, y colaboradores de las partes como el patrocinio letrado, el apoderamiento (arts. 47 CPCCN y 1017 CCCN) y los consultores técnicos.

Sobre esa base se construyen los ejes de competencia territorial (contrato, prórroga tácita, reglas generales de los arts. 5 y siguientes del CPCCN) y de calidad y legitimidad de parte, distinguiendo con precisión capacidad de hecho, capacidad de derecho, representación legal y representación convencional. Un bloque completo se dedica a los distintos domicilios que conviven en un mismo proceso --- especial, procesal o \emph{ad litem}, legal, electrónico y real --- y a las consecuencias prácticas de confundirlos.

La materia repasa también los principios que estructuran el trámite: cargas, deberes y obligaciones, principio dispositivo, sistema de preclusión, litisconsorcio (voluntario y necesario) y una mención acotada al proceso colectivo, útil para diferenciarlo de los desalojos administrativos y sus límites frente a la emergencia habitacional.

Cierra con dos herramientas de fuerte aplicación práctica: la extinción de la fianza por falta de renovación formal del contrato (art. 1225 CCCN) y el principio de disponibilidad de las formas, que permite a las partes acordar mecanismos alternativos --- como una pericial arbitral --- para resolver puntos controvertidos sin necesidad de agotar toda la instancia probatoria ordinaria.

El criterio transversal que atraviesa todo el desarrollo es que el conocimiento técnico solo tiene valor si se sabe aplicar oportunamente.

\end{document}
